% !TeX program = pdflatex
\documentclass[11pt,a4paper]{article}
\usepackage{../../recursos/latex/estilo-notas}

\begin{document}
\cabecalhoaula{10}{KNN e Árvores de Classificação}{AME §8.3--8.6; ISLP §4.7.6 e cap.~8}

%==============================================================================
\section*{Objetivos da aula}
%==============================================================================
Ao final desta aula o aluno deve ser capaz de:
\begin{itemize}[leftmargin=1.4cm]
  \item adaptar o KNN para classificação, pelo voto da maioria, e estimar
        $\Prob(Y=c\mid\x)$ pelas proporções das classes entre os vizinhos;
  \item construir uma árvore de classificação e justificar a predição pela moda;
  \item definir as medidas de impureza de Gini e de entropia e justificar por que
        são preferidas à taxa de erro;
  \item usar florestas e \emph{boosting} em classificação (voto, $m$, AdaBoost);
  \item comparar os classificadores vistos no curso quanto à fronteira, à
        interpretabilidade e ao desempenho.
\end{itemize}

\begin{emsala}
A aula fecha a classificação reaproveitando dois métodos não paramétricos vistos em
regressão, o KNN (Aula~04) e as árvores (Aula~05). A estrutura dos dois não muda; muda
o resumo de cada vizinhança ou folha, em que a média dá lugar à moda, e o critério de
divisão das árvores, em que o EQM dá lugar à impureza.
\end{emsala}

%==============================================================================
\section{KNN para classificação}
%==============================================================================
Seja $\mathcal{N}_{\x}^{(k)}=\{i : d(\X_i,\x)\le d^k_{\x}\}$ o conjunto dos índices das
$k$ observações de treino mais próximas de $\x$, como na Aula~04. O classificador dos
$k$ vizinhos mais próximos atribui a $\x$ a classe mais frequente entre os vizinhos,
pelo voto da maioria:
\begin{equation}\label{eq:knncls}
  \ghat(\x) = \operatorname{moda}\{Y_i : i\in\mathcal{N}_{\x}^{(k)}\}.
\end{equation}
É um classificador \emph{plug-in} (Aula~07). A estimativa
\[
  \widehat\Prob(Y=c\mid\x) = \frac{1}{k}\sum_{i\in\mathcal{N}_{\x}^{(k)}}\1{Y_i=c},
\]
a proporção de vizinhos da classe $c$, está em $[0,1]$ e soma $1$ sobre as classes;
como $k\,\widehat\Prob(Y=c\mid\x)$ é o número de votos da classe $c$, a classe de maior
probabilidade estimada é a mais votada. E, como
$\Prob(Y=c\mid\X=\x)=\E\big[\1{Y=c}\mid\X=\x\big]$, essa estimativa é o KNN de regressão
da Aula~04 aplicado à indicadora $\1{Y=c}$.

\begin{ideia}
$k$ controla o balanço entre viés e variância, como na regressão: $k=1$ dá uma
fronteira muito recortada, com viés baixo e variância alta; $k$ grande dá uma
fronteira suave, com viés alto e variância baixa. Escolha $k$ por validação cruzada
(Aula~03).
\end{ideia}

\begin{atencao}
Valem os cuidados da Aula~04: \emph{padronize} as covariáveis, porque o KNN se baseia
numa distância, e lembre da \emph{maldição da dimensionalidade}: o KNN não seleciona
covariáveis e piora com muitas covariáveis irrelevantes. Em problemas binários, use
$k$ ímpar para evitar empates na votação.
\end{atencao}

%==============================================================================
\section{Árvores de classificação}
%==============================================================================
São o análogo das árvores de regressão (Aula~05): uma partição recursiva do espaço das
covariáveis em regiões disjuntas $R_1,\dots,R_J$. Para uma região $R$, sejam $n_R$ o
número de observações de treino em $R$ e $\widehat p_{R,c}$ a proporção da classe $c$
entre elas. Há duas adaptações.

\subsection{Predição: a moda}
Numa folha $R_k$, prediz-se a classe majoritária, em vez da média:
\begin{equation}\label{eq:arvcls}
  \ghat(\x) = \operatorname{moda}\{Y_i : \X_i\in R_k\},\qquad \x\in R_k,
\end{equation}
e $\widehat\Prob(Y=c\mid\x)=\widehat p_{R_k,c}$, a proporção da classe $c$ na folha.

\subsection{Critério de divisão: impureza}
\begin{emsala}
Pergunta para a turma, antes das fórmulas: \emph{``como uma árvore decide a melhor
pergunta a fazer em cada nó quando o alvo é uma classe?''}
\end{emsala}

\noindent O EQM não faz sentido com $Y$ qualitativo. Buscam-se divisões que tornem as
regiões \emph{puras}, idealmente com uma só classe, e mede-se a impureza de $R$ por
\begin{align}
  \text{Índice de Gini:}\quad & G(R) = \sum_{c\in\mathcal{C}} \widehat p_{R,c}\,(1-\widehat p_{R,c}),
  \label{eq:gini}\\
  \text{Entropia:}\quad & D(R) = -\sum_{c\in\mathcal{C}} \widehat p_{R,c}\,\log_2 \widehat p_{R,c},
  \label{eq:entropia}
\end{align}
com a convenção $0\log_2 0=0$. As duas se anulam se e somente se $R$ é pura (algum
$\widehat p_{R,c}=1$) e são máximas quando as classes estão equilibradas, com
$G(R)\le1-1/\abs{\mathcal{C}}$ e $D(R)\le\log_2\abs{\mathcal{C}}$. Com $I$ igual a $G$
ou a $D$, a divisão de $R$ em $R_1$ e $R_2$, com $n_1$ e $n_2$ observações, é escolhida
de modo a maximizar a redução de impureza
\[
  \Delta I = I(R)-\Big[\frac{n_1}{n_R}\,I(R_1)+\frac{n_2}{n_R}\,I(R_2)\Big],
\]
o que equivale a minimizar a impureza ponderada $n_1I(R_1)+n_2I(R_2)$.

\begin{observacao}[Por que não usar a taxa de erro?]
A taxa de erro de classificação $E(R)=1-\max_c \widehat p_{R,c}$ também mede
impureza, mas é insensível a mudanças nas proporções que não alterem a classe
majoritária: no caso binário, ela é linear em cada metade de $[0,1]$, e, se os dois
filhos têm a mesma classe majoritária, a média ponderada dos erros deles é igual ao
erro do pai. Gini e entropia são estritamente côncavas e respondem a qualquer melhora
na pureza; por isso fazem divisões melhores. A taxa de erro é boa para a \emph{poda}
final.
\end{observacao}

\begin{exemplo}[Escolha de divisão pelo Gini]
Para predizer se um hospital tem tomógrafo, o [AME] compara, em dez hospitais, a
divisão por especialidade em cardiologia com a divisão por mais de $100$ leitos. O
Gini ponderado da divisão por cardiologia é
$\frac4{10}\cdot\frac12+\frac6{10}\cdot\frac49=\frac7{15}\approx0{,}467$, e o da
divisão por leitos é
$\frac6{10}\cdot\frac5{18}+\frac4{10}\cdot\frac38=\frac{19}{60}\approx0{,}317$; a de
\emph{menor} valor, a mais pura, é a divisão por leitos. O [AME] dá $0{,}472$ e
$0{,}326$, porque soma $\widehat p(1-\widehat p)$ nas duas regiões sem ponderar pelo
tamanho delas; a escolha é a mesma.
\end{exemplo}

\noindent A construção segue as duas etapas da Aula~05: \emph{cresce-se} uma árvore
grande e depois ela é \emph{podada} por custo--complexidade, com $\alpha$ escolhido por
validação cruzada. Valem as mesmas vantagens: interpretabilidade, tratamento direto
de covariáveis categóricas, seleção automática de covariáveis e captura de
interações.

%==============================================================================
\section{\emph{Ensembles} para classificação}
%==============================================================================
As formas de agregar árvores valem também em classificação, agora com árvores de
classificação:
\begin{description}[leftmargin=2.6cm,style=nextline]
  \item[Bagging / Florestas] $B$ árvores ajustadas em amostras \emph{bootstrap},
        agregadas pelo voto da maioria ou pela média das probabilidades estimadas.
        As florestas diminuem a correlação entre as árvores sorteando $m$
        covariáveis candidatas em cada nó; em classificação, a regra empírica é
        $m\approx\sqrt{p}$. Aumentar $B$ não leva a superajuste, e a estimativa OOB
        do risco e a importância das covariáveis saem sem custo adicional.
  \item[Boosting] árvores ajustadas em sequência. A variante clássica em
        classificação é o AdaBoost: a cada rodada, as observações são reponderadas,
        com mais peso para as mal classificadas, e o classificador final é um voto
        ponderado dos aprendizes fracos. Um $B$ grande pode levar a superajuste (use
        parada antecipada, \emph{early stopping}); o XGBoost é a implementação de
        referência.
\end{description}

%==============================================================================
\section{Comparando os classificadores do curso}
%==============================================================================
\begin{center}\small
\renewcommand{\arraystretch}{1.3}
\begin{tabular}{@{}l>{\raggedright\arraybackslash}p{3.5cm}>{\raggedright\arraybackslash}p{3.2cm}>{\raggedright\arraybackslash}p{3.6cm}@{}}
\toprule
\textbf{Método} & \textbf{Fronteira} & \textbf{Interpretab.} & \textbf{Observações} \\
\midrule
Logística & linear (no \emph{log-odds}) & alta (coefs.) & dá probabilidades; robusta \\
LDA / QDA & linear / quadrática & média & ótimos sob gaussianidade \\
Bayes ingênuo & depende & média & forte com texto; indep.\ condicional \\
KNN & muito flexível, local & baixa & sensível a escala e dimensão \\
Árvore & ``escada'' (eixos) & \textbf{altíssima} & instável sozinha \\
Florestas / Boosting & muito flexível & baixa & topo de desempenho em dados tabulares \\
SVM (kernel) & linear ou não linear & baixa & ótima com classes separadas \\
\bottomrule
\end{tabular}
\end{center}

\begin{ideia}
Nenhum método é melhor que os outros em todos os problemas (``não existe almoço
grátis''). A escolha depende da estrutura dos dados e do objetivo: são necessárias
probabilidades calibradas? interpretabilidade? há muitas covariáveis irrelevantes? a
fronteira é não linear? Compare alguns candidatos por validação cruzada, com as
métricas adequadas (Aula~08).
\end{ideia}

\begin{emsala}
\raggedright
Use os notebooks \texttt{Aula prática 10}, \texttt{Comparação entre classificadores
paramétricos} e \texttt{Exemplo - KNN e árvores em dados artificiais} para mostrar as
fronteiras de decisão lado a lado. É o melhor jeito de fixar a oposição entre métodos
flexíveis e rígidos, que percorreu o curso inteiro.
\end{emsala}

\begin{emsala}[o item ``probabilidades calibradas'' não é preferência]
A §7 da \texttt{Aula prática 10} refaz a tabela numa base em que a classe positiva é
rara ($10\%$ das observações). Vale mostrar as duas medições:

\begin{itemize}[leftmargin=1.2cm,nosep]
  \item a acurácia varia de $0{,}8989$ a $0{,}9072$ entre os sete métodos, e quatro
        deles empatam exatamente com quem responde ``não'' para todos, porque nunca
        preveem a classe positiva; na mesma tabela, a AP varia de $0{,}119$ a
        $0{,}453$;
  \item com o corte de custo $1{:}20$ da Aula~08, $1/21\approx0{,}0476$, a árvore
        podada e o AdaBoost mandam investigar todos os clientes, porque as
        probabilidades deles quase nunca ficam abaixo de $0{,}05$; a cada mil
        clientes, entre os quais há $101$ positivos, a logística manda investigar
        $410$ e encontra $82$, e a SVM manda $493$ e encontra $89$.
\end{itemize}

O ponto a registrar no quadro: o corte ótimo é ótimo para probabilidades
calibradas. Pergunte à turma o que fazer com um modelo que ordena bem e calibra mal;
a resposta vem da Aula~08, e é calibrar, não trocar de modelo.
\end{emsala}

%==============================================================================
\section*{Resumo}
%==============================================================================
\begin{itemize}[leftmargin=1.4cm]
  \item KNN para classificação \eqref{eq:knncls}: voto da maioria; estima
        $\Prob(Y=c\mid\x)$ pelas proporções entre os vizinhos; $k$ por validação
        cruzada; padronize as covariáveis.
  \item Árvore de classificação: prediz a moda \eqref{eq:arvcls}; divide pelo Gini
        \eqref{eq:gini} ou pela entropia \eqref{eq:entropia}, melhores que a taxa de
        erro; cresce e é podada como na Aula~05.
  \item Florestas ($m\approx\sqrt p$) e \emph{boosting} (AdaBoost) agregam árvores por
        voto e estão entre os métodos de melhor desempenho em dados tabulares.
  \item Nenhum classificador é melhor em todos os problemas: escolha pela estrutura
        dos dados e pelo objetivo, e compare por validação cruzada com as métricas da
        Aula~08.
\end{itemize}

\paragraph{Leitura recomendada.} [AME] §8.3 (árvores de classificação, índice de
Gini), §8.4 e §8.5 (\emph{bagging}, florestas, \emph{boosting}), §8.6 (KNN); [ISLP]
§4.7.6 (KNN) e Capítulo~8 (árvores e \emph{ensembles}, incluindo critérios de
impureza e \emph{boosting}).

\end{document}
